\section*{Part V: Hiring Consultants, Operators, and Salespeople}

\subsection*{45. A Consultant Is Selected for Problem Ownership, Not Slide Production}

A consultant may help a client understand a problem, make a decision, execute a technical intervention, or develop internal capability. The job depends on \textbf{diagnosis, analytical rigor, contextual judgment, stakeholder communication, delivery discipline, and ethical scope control}. A candidate who speaks in elegant frameworks but cannot distinguish an assumption from a fact creates confident confusion. A candidate who can technically solve the problem but ignores the buyer's constraints may deliver an unusable answer.

The assessment should recreate a client's ambiguity: a business objective, incomplete data, conflicting stakeholder incentives, and a deadline. Ask the candidate to decide what must be learned before recommending anything, what can be done with current evidence, and what should be explicitly excluded. For an AI transformation consultant, do not let the performance end with a chatbot demo; require workflow baselines, security checks, user adoption, and whether the proposed improvement is worth its total operating cost.

\subsection*{46. A 30-Minute Consulting Discovery Simulation}

Provide a fictional client brief: ``Our support team is too slow. We want an AI agent.'' The candidate interviews a trained role-player who knows that most delays arise from exceptions requiring supervisor approval, not from writing responses. The candidate has 12 minutes to discover the operational problem, 8 minutes to structure hypotheses, and 10 minutes to recommend a next step.

A strong candidate tests the alleged bottleneck, asks for volume and variance, identifies who owns decisions, distinguishes automatable from regulated or sensitive work, estimates the value of faster handling, and avoids promising a technology-led solution prematurely. A weak candidate pivots immediately to vendor features and lists fashionable use cases. Score quality of questions, diagnostic decomposition, risk awareness, economic framing, and ability to say ``we do not yet know.'' This is also a test of intellectual honesty under client pressure.

\subsection*{47. Give Consultants a Deliverable with an Uncomfortable Conclusion}

An effective consultant must sometimes tell a client that the proposed investment should not proceed. Provide evidence for an initiative with a convincing executive narrative but poor counterfactual economics. Ask for a 1-page recommendation, an assumptions register, and a suggested validation experiment. See whether the candidate can respectfully recommend deferral or discontinuation, while preserving a constructive client relationship.

This is where sales ability and advisory independence can conflict. A consultant paid for implementation may be biased toward an implementation recommendation. Recruitment should therefore test how the individual identifies and manages conflicts of interest. Ask for a real incident where they recommended against the most lucrative next step. Probe for the decision, stakeholders, alternatives, and consequences: not merely a statement of values.

\subsection*{48. The Consultant Scorecard}

A suggested framework is:
\begin{itemize}[leftmargin=*]
\item Problem framing: 25\%
\item Analytical and technical correctness: 25\%
\item Actionable recommendation: 15\%
\item Stakeholder communication: 15\%
\item Delivery discipline: 10\%
\item Ethical independence: 10\%
\end{itemize}
with must-pass standards for truthfulness, client confidentiality, and the essential domain knowledge required at entry.

For a technical consultant, require an implementation or code artifact; for an organizational consultant, require evidence of changing behavior and governance; for an education consultant, assess learning outcomes and institutional constraints. The candidate's presentation can be scored, but it should not overwhelm the substantive quality of the recommendation. References should address whether the candidate delivered what was agreed, handled uncertainty responsibly, and left the client more capable rather than more dependent.

\subsection*{49. Operators Are Measured When the Normal Plan Fails}

An operations specialist transforms a service promise into reliable execution. The work often looks ordinary until a customer escalation, delayed vendor, instructor absence, failed payment, conflicting schedule, or data error threatens delivery. Important competencies include prioritization, attention to detail, process control, communication, escalation, quality assurance, and recovery after failure.

A generic ``describe a time you worked under pressure'' question is not enough. Give a simulated operating day with a queue of incoming requests and finite resources. Let the candidate construct a plan, discover a new incident halfway through, and revise the plan. Observe whether they protect safety and critical customers, use timestamps and ownership, communicate realistic expectations, and avoid hiding bad news. Good operations work is partly the ability to notice the exception before it becomes a crisis.

\subsection*{50. An Operations In-Tray That Reveals Priority Judgment}

Imagine a hypothetical learning service at 9:00~a.m.: a live session starts in 2 hours; an instructor reports illness; 3 learners cannot access assignments; a promised refund is overdue; an enterprise buyer requests a curriculum change; the founder is unavailable until evening. The candidate gets a 1-page resource and authority map. Ask them to write the 1st 10 actions and messages, including what should \textbf{not} be done.

There is no unique perfect sequence. But the evaluation should reward identification of hard deadlines, learner harm, contractual obligations, alternatives, delegated owners, and escalation triggers. A weak answer promises every stakeholder immediate resolution. A strong 1 communicates realistic recovery, keeps a record, chooses a substitute instructor or approved contingency, and protects the integrity of assessments. Time-stamped decisions are better evidence than a confident claim of being ``organized.''

\subsection*{51. Measure Process Improvement Separately from Firefighting}

A heroic operator may save every crisis while failing to remove recurring causes. A more valuable operator designs incident categories, service-level definitions, checklists, monitoring, staff backups, root-cause reviews, and clear decision rights. Ask for a 30-day improvement plan based on a fictional error log. Candidates should prioritize preventable recurring incidents and distinguish failures of process from failures of individual effort.

For a very junior role, reliable execution of documented tasks may be the primary threshold. For an operations lead, the ability to \textbf{make average performance better without constant personal intervention} matters more. Ask how the candidate would train a replacement, monitor error rates, and know when additional staff are needed. ``I would personally make sure it never happens again'' is often evidence of a future bottleneck.

\subsection*{52. The Operator Scorecard}

An illustrative service-operations rubric is:
\begin{itemize}[leftmargin=*]
\item Prioritization and exception handling: 25\%
\item Accuracy and procedural control: 20\%
\item Customer communication: 20\%
\item Documentation and handoffs: 15\%
\item Process improvement: 10\%
\item Reliability and ownership: 10\%
\end{itemize}
For payment processing, privacy, or safety tasks, define non-negotiable compliance gates.

Assess both a routine artifact (a reconciled roster, schedule, or ticket batch) and an exception exercise. A candidate might be accurate when unhurried but unable to adapt under uncertain conditions; another may be excellent in crises but careless with repeatable controls. Match the mix to the role. For early Bayes, a part-time coordinator is valuable only if their work can be bounded and they can be entrusted with personal data, learner communication, and reliable follow-through.
